\section{Symmetric partner averaging and the exponent $2^{-145}$}
\label{sec:two-scale-symmetric-pairs}

The far-pair estimate improves when it is restricted to the symmetric family of multiple-link pairs.
This reduces the initial cycle mass needed to produce a chain from two to one.
The resulting parameter choice gives a larger explicit exponent.

\begin{lemma}[Symmetric partner averaging]
\label{lem:symmetric-partner-averaging}
In either cycle family of the two-scale argument, let $M$ be the total weight of ordered disjoint pairs
with no edge between them and at least two links.
Suppose the local intersection and direct-edge weights are $o(\mu)$, and the one-link weight is at most $(1+o(1))\mu$.
Then some cycle has main-component multiple-link partners of total weight at least
\[
 \frac{\mu-1}{2}-o(1).
\]
The main component is defined relative to one fixed connected marked component.
\end{lemma}

\begin{proof}
The hypotheses give $M\ge\mu^2-(1+o(1))\mu$.
The relation defining $M$ is symmetric.
Within each unordered pair, at most one orientation can be far, by the far-pair antisymmetry argument.
Thus far orientations have weight at most $M/2$, rather than merely $\mu^2/2$.
The main orientations have weight at least
\[
 M/2\ge\frac12\mu^2-\frac{1+o(1)}2\mu.
\]
Divide by the total first-cycle weight $\mu$ to obtain the conclusion.
\end{proof}

\begin{theorem}[A revised explicit exponent]
\label{thm:two-scale-exponent-145}
Put $\widehat\kappa=2^{-145}$.
Under the degree hypotheses of Theorem~\ref{thm:two-scale-critical-marking}, suppose
\[
 |P|\le\frac{r}{15000\log_2r},\qquad
 \operatorname{comp}(\Gamma[P])\le\frac{\log_2r}{1000}.
\]
For all sufficiently large $r$,
\[
 \Pr(\mathcal F)\le\frac12+r^{-\widehat\kappa}.
\]
Consequently,
\[
 \psi(r)\le2^{r-1}+(1+o(1))\frac{2^r}{r^{\widehat\kappa}}.
\]
\end{theorem}

\begin{proof}
Run the same two-scale proof with
\[
 A=15,\qquad t=128,\qquad x=\log_2r,\qquad L=\lfloor15x\rfloor,
\]
\[
 D_0=\lfloor49x/100\rfloor,\qquad
 D_1=\lfloor99x/100\rfloor,
 \qquad s_j=\lfloor(D_j-1)/2\rfloor.
\]
Assume that the forest probability exceeds $1/2+r^{-\widehat\kappa}$.
The nested deletions satisfy
\[
 |S_0|=O(x^3r^{49/100}),\qquad
 |S_1|=O(x^3r^{99/100}).
\]

Use paired spectral powers starting at the smallest even integer at least $2x+10\log_2x$.
End at the largest odd integer at most $L$.
The marked-set bound again gives $Ld_1/(2v_1)<1/500$ for large $r$.
The nonreal error is $O(x^{-6})$, and the repeated-walk factor is $1+o(1)$.
Therefore
\[
 \mu_1\ge\frac{499}{500}\,\frac{\ln(15/2)}2-o(1)
 >\frac{201}{200}=1.005.
\]
The limiting numerical lower bound is $1.005436607250590\ldots$.

At both girth scales, the link decomposition has asymptotic coefficient at most
\[
 \frac1{1000}+\frac{30}{128}
       +8\cdot128\cdot2^{128}\widehat\kappa\ln2
 =\frac1{1000}+\frac{30}{128}+\frac{\ln2}{128}
 <\frac{61}{250}=0.244.
\]
The limiting value is $0.240790212348124\ldots$.
The bounded tree-link term and all rounding errors are absorbed for large $r$.
Since $s_0=49x/200+O(1)$, exceptional partners contribute
\[
 O(a_Cx^3r^{-1/1000})=o(a_C).
\]
Both inverse moment estimates remain valid.

Lemma~\ref{lem:symmetric-partner-averaging} now supplies a cycle with main-component partner mass at least $1/1000$.
It yields a nested chain of length
\[
 T=\Omega(r^{99/200}/x^3).
\]
Grouping this chain into full blocks gives $\Omega(r^{99/200}/x^4)$ candidate blocks.
The discarded marked components, anchor neighbors, and $S_0$ neighborhoods total
\[
 O(x^3r^{49/100})=o(r^{99/200}/x^4).
\]
Thus the clean-block construction works unchanged.

For the multiple-mode spectral test, keep $\theta=1/(12x)$ but use even lengths between $6x$ and $15x$.
Its lower-bound coefficient per real eigenvalue tends to
\[
 e^{-15/24}\frac{\ln(15/6)}4
 =0.122613771520417\ldots>\frac1{16}.
\]
The nonreal and repeated-walk errors still vanish.
Hence
\[
 \mu_0=\Omega(r^{99/200}/x^4).
\]
The inverse moment bound gives $\mu_0\le(1+o(1))r^{\widehat\kappa}/4$, a contradiction.
The direct cubic-extremizer transfer proves the cycle estimate.
\end{proof}

These constants remain too large for a practical finite-rank bound.
The improvement changes the explicit exponent and the marked-set coefficient, while preserving the underlying proof.
